\documentclass[11pt,letterpaper]{article}
\usepackage{technical-user-manual}
\setDocTitle{Build Reproducibility, Dependency and Cache Engineering}
\setDocSubtitle{Engineering technical reference}
\setDocVersion{1.0}
\setDocStatus{Technical education / proposed reference}
\setDocOwner{CI/CD Engineering}
\setDocDate{30 October 2026}
\setDocSummary{Supplementary engineering practice}
\begin{document}
\makedoctitle
\makedocsummary
\tableofcontents
\clearpage
\section{Purpose, scope and authority}
This is an original modern engineering reference. The 2007 Duvall, Matyas and Glover book informs foundational CI intent but does not prescribe GitHub-specific implementations or 2026 supply-chain standards.

\section{Reproducibility levels}
Repeatable builds use controlled commands; hermetic builds additionally limit undeclared network and environment inputs; byte-for-byte reproducibility demands deterministic timestamps, archive ordering, toolchains and dependencies. Do not label ordinary caches as hermeticity.

Verification: collect the configuration or runtime evidence supporting this control, identify an accountable owner, and record the exception path where the condition cannot presently be met.

\section{Dependency policy}
Commit applicable lockfiles and verify resolved versions and checksums. Use an approved internal proxy if needed. Track transitive dependency changes and intentional exceptions. Never auto-upgrade or silently disable audit gates during a release.

Verification: collect the configuration or runtime evidence supporting this control, identify an accountable owner, and record the exception path where the condition cannot presently be met.

\section{Cache threat model}
Caches trade latency for reuse and are not an authority for correctness. Trust isolation matters because restored objects can be poisoned; never cache access tokens or secrets. Treat cache key, scopes, branch trust and restore prefixes as reviewable design choices.

Verification: collect the configuration or runtime evidence supporting this control, identify an accountable owner, and record the exception path where the condition cannot presently be met.

\section{Pipeline stage allocation}
Stage A performs formatting, policy checks and static validation. Stage B executes unit tests and language builds. Stage C performs integration, quality and security checks. Stage D packages and signs/attests a candidate. Slow checks may be asynchronous only if protection rules prevent bypass.

Verification: collect the configuration or runtime evidence supporting this control, identify an accountable owner, and record the exception path where the condition cannot presently be met.

\section{Measurement protocol}
Measure cold/warm cache separately, runners and toolchain versions fixed, median and tail runtime, download size, CPU/memory constraints, failures and misses. Compare paired baseline/candidate conditions with declared seeds; avoid treating single runs as statistically meaningful.

Verification: collect the configuration or runtime evidence supporting this control, identify an accountable owner, and record the exception path where the condition cannot presently be met.

\section{Adoption checklist}
\begin{itemize}
\item Identify applicable repositories, application owners and trust boundaries.
\item Define minimum pass criteria and measurable acceptance thresholds.
\item Review workflow permissions, external actions, outputs and retention.
\item Record exceptions with justification, owner and expiry.
\item Revalidate after action, runner or policy changes.
\end{itemize}
\end{document}
